\documentclass[12pt,a4paper]{amsart}

\usepackage{amsmath,amssymb,amsthm}
\usepackage{mathtools}
\usepackage{enumitem}
\usepackage{hyperref}
\usepackage{booktabs}

%%%─── Theorem environments
\newtheorem{theorem}{Theorem}[section]
\newtheorem{lemma}[theorem]{Lemma}
\newtheorem{proposition}[theorem]{Proposition}
\newtheorem{corollary}[theorem]{Corollary}
\newtheorem{conjecture}[theorem]{Conjecture}
\theoremstyle{definition}
\newtheorem{definition}[theorem]{Definition}
\newtheorem{assumption}[theorem]{Assumption}
\newtheorem{problem}[theorem]{Open Problem}
\theoremstyle{remark}
\newtheorem{remark}[theorem]{Remark}

%%%─── Macros
\newcommand{\Kc}{\mathcal{K}_c}
\newcommand{\Tr}{\mathrm{Tr}}
\newcommand{\lam}[2]{\lambda_{#2}^{(#1)}}
\newcommand{\gap}[2]{\mathrm{gap}_{#2}^{(#1)}}
\newcommand{\Ai}{\mathrm{Ai}}
\newcommand{\op}{\mathrm{op}}
\newcommand{\C}{\mathbb{C}}
\newcommand{\R}{\mathbb{R}}

\subjclass[2020]{47B34, 45C05, 35P20, 47A55, 34E20, 81Q20}
\keywords{Prolate spheroidal wave functions, sinc kernel, resolvent asymptotics,
  Airy operator, double-scaling limit, Fredholm determinant, spectral counting}

\begin{document}

\title[Paper XIV: Airy Resolvent for Sinc Kernel]{%
  Paper~XIV: Airy Resolvent Asymptotics for the Sinc Kernel Operator\\[0.3em]
  \large Towards a Proof of the Local Spectral Counting Condition~(Q5)}

\author{\textit{prolate-gram-coercivity programme}}
\date{May 2026}

\begin{abstract}
Paper~XIII identified condition~(Q5) --- the local spectral counting
statement $N_\varphi(c) = 2$ in a $c^{-1/3}$-scale transition window ---
as the primary obstruction to Gap-S.
In its binary form, (Q5) is a reformulation of the spectral problem,
not a reduction from independent principles.
This paper initiates the programme to convert~(Q5) into a genuine
reduction by controlling
\[
  \Tr[\varphi_c(\Kc)] =
  -\frac{1}{\pi}\int_{\C}\partial_{\bar z}\tilde\varphi_c(z)\,
  \Tr[(\Kc - z)^{-1}]\,d^2z
\]
directly from the integral-operator structure of~$\Kc$.

\medskip\noindent
\textbf{Central target.}
Derive the leading asymptotic expansion
\[
  \Tr[(\Kc - z)^{-1}] \sim
  c\cdot R_0(z) + c^{2/3}\cdot R_{\mathrm{Ai}}(z;\,\xi) + O(c^{1/3})
\]
in the double-scaling regime
$z = z_n + c^{-1/3}\zeta$, $z_n = \lam{c}{n}$,
where $R_{\mathrm{Ai}}$ is the resolvent of the Airy operator
$\mathcal{A} = -d^2/d\xi^2 + \xi$ on $L^2(\R_+)$.

\medskip\noindent
\textbf{If achieved}, this gives:
\begin{itemize}
  \item $N_\varphi(c) = \Tr[\varphi_c(\Kc)] = 2 + O(c^{-1/3})$
    for cutoffs $\varphi_c$ adapted to the $c^{-1/3}$-window;
  \item condition~(Q5) from integral-operator asymptotics alone,
    independent of the eigenvalues themselves;
  \item Gap-S under the remaining assumptions of Theorem~B (Paper~XIII).
\end{itemize}
\end{abstract}

\maketitle
\tableofcontents
\bigskip\hrule\bigskip

%% ============================================================
\section{Setup and Goal}
\label{sec:setup}
%% ============================================================

\subsection{The operator and the transition regime}

The prolate concentration operator is
\begin{equation}\label{eq:Kc}
  (\Kc f)(x) = \int_{-1}^1 \frac{\sin(c(x-y))}{\pi(x-y)} f(y)\,dy,
  \quad f \in L^2([-1,1]),
\end{equation}
with eigenvalues $1 > \lam{c}{0} \geq \lam{c}{1} \geq \cdots > 0$
and the Weyl-type count
\begin{equation}\label{eq:Weyl}
  \#\{n : \lam{c}{n} > \mu\} \sim \frac{2c}{\pi}\quad (\mu \approx 1/2,\;c\to\infty).
\end{equation}
The \emph{transition index} satisfies $n^* \sim 2c/\pi$;
the \emph{transition regime} is $n = n^* + k$ with $k \sim c^{1/3}$.

\subsection{The missing theorem}

Define the \emph{transition window} centred at $z_n^* := \lam{c}{n^*}$:
\begin{equation}\label{eq:window}
  W_c(\zeta_0) :=
  \bigl[z_n^* - \zeta_0 c^{-1/3},\; z_n^* + \zeta_0 c^{-1/3}\bigr],
  \quad \zeta_0 > 0 \text{ fixed}.
\end{equation}
For a smooth cutoff $\varphi_c \in C_c^\infty(W_c(\zeta_0))$, $0 \leq \varphi_c \leq 1$,
equal to $1$ on $W_c(\zeta_0/2)$, define
\begin{equation}\label{eq:smoothcount}
  N_\varphi(c) := \Tr[\varphi_c(\Kc)] = \sum_j \varphi_c(\lam{c}{j}).
\end{equation}

\begin{problem}[Main target of Paper~XIV]
\label{prob:main}
Prove
\begin{equation}\label{eq:target}
  N_\varphi(c) = 2 + O(c^{-1/3})
\end{equation}
from the integral-operator structure of $\Kc$,
without reference to the individual eigenvalues $\lam{c}{j}$.
\end{problem}

This is condition~(Q5) of Paper~XIII~\cite{paper13},
now reformulated as a \emph{trace asymptotics problem}.

%% ============================================================
\section{The Double-Scaling Limit}
\label{sec:doublescaling}
%% ============================================================

\subsection{Scaling ansatz}

In the transition regime, the natural variables are:
\begin{equation}\label{eq:scaling}
  x = 1 - c^{-2/3}\xi, \quad
  y = 1 - c^{-2/3}\eta, \quad
  z = z_n^* + c^{-1/3}\zeta,
\end{equation}
where $\xi,\eta \in \R_+$ and $\zeta \in \C\setminus\R$.
This is the \emph{Airy double-scaling}: the spatial variable
scales as $c^{-2/3}$ (same as classical turning-point scale),
and the spectral variable scales as $c^{-1/3}$.

\begin{remark}[Motivation from random matrix theory]
For the GUE and the sine kernel, an analogous double-scaling at the
spectral edge produces the Airy kernel operator
\begin{equation}\label{eq:Airykernel}
  K_{\Ai}(\xi,\eta) =
  \frac{\Ai(\xi)\Ai'(\eta) - \Ai'(\xi)\Ai(\eta)}{\xi - \eta}
\end{equation}
as the universal limiting kernel (Tracy--Widom).
For the sinc kernel $K_c(x,y) = \sin(c(x-y))/(\pi(x-y))$,
a similar but distinct double-scaling at the endpoint $x=\pm 1$
produces (heuristically) a modified Airy-type operator.
The key difference: the sinc kernel is not translation-invariant
on $[-1,1]$, so universality does not follow automatically.
\end{remark}

\subsection{The target resolvent expansion}

Let $R(z;c) := (\Kc - z)^{-1}$ for $\mathrm{Im}(z) \neq 0$.
In the double-scaling limit~\eqref{eq:scaling}, we expect:
\begin{equation}\label{eq:resolventexp}
  \boxed{
  R(z_n^* + c^{-1/3}\zeta;\, c)
  \;\sim\;
  c\cdot R_0(z_n^*)
  \;+\;
  c^{2/3}\cdot R_{\mathrm{Ai}}(\zeta)
  \;+\;
  c^{1/3}\cdot R_1(\zeta)
  \;+\; O(1)
  }
\end{equation}
in an appropriate operator norm, where:
\begin{enumerate}[label=(\roman*)]
  \item $R_0(z_n^*)$: the \emph{bulk resolvent}, capturing the
    global eigenvalue density away from the transition;
  \item $R_{\mathrm{Ai}}(\zeta) = (\mathcal{A} - \zeta)^{-1}$:
    the \emph{Airy resolvent}, where
    $\mathcal{A} = -d^2/d\xi^2 + \xi$ on $L^2(\R_+)$
    with Dirichlet boundary condition at $\xi=0$;
  \item $R_1(\zeta)$: a correction term depending on the
    specific geometry of the interval $[-1,1]$;
    this term is absent for the sine kernel on $\R$.
\end{enumerate}

\begin{remark}[Status of \eqref{eq:resolventexp}]
Expansion~\eqref{eq:resolventexp} is a \emph{conjecture}.
It is motivated by:
\begin{itemize}[nosep]
  \item The Tracy--Widom scaling for the GUE hard edge;
  \item Semiclassical turning-point theory for the differential
    equation satisfied by PSWFs
    (prolate spheroidal wave equation, PSWE);
  \item The Widom-type formula for the global trace:
    $\Tr[\Kc] \sim 2c/\pi + O(\log c)$ (Widom~\cite{Widom1964}).
\end{itemize}
Verifying~\eqref{eq:resolventexp} rigorously is the central analytic
goal of this paper.
\end{remark}

\subsection{Trace asymptotics from the resolvent expansion}

Assume~\eqref{eq:resolventexp} holds in trace norm.
Applying the Helffer--Sj\"ostrand formula
\begin{equation}\label{eq:HS}
  N_\varphi(c) =
  -\frac{1}{\pi}\int_{\C}\partial_{\bar z}\tilde\varphi_c(z)\,
  \Tr[R(z;c)]\,d^2z
\end{equation}
with $\tilde\varphi_c$ an almost-analytic extension of $\varphi_c$
supported in a $c^{-1/3}$-neighbourhood of $W_c$, and
changing variables $z = z_n^* + c^{-1/3}\zeta$:
\begin{equation}\label{eq:tracescaled}
  N_\varphi(c) =
  -\frac{c^{-1/3}}{\pi}\int_{\C}
  \partial_{\bar\zeta}\tilde\varphi(\zeta)\,
  \Tr[R(z_n^*+c^{-1/3}\zeta;c)]\,d^2\zeta.
\end{equation}
Inserting~\eqref{eq:resolventexp}:
\begin{equation}\label{eq:traceexpand}
  N_\varphi(c) \sim
  \underbrace{-\frac{c^{2/3}}{\pi}\int_\C
    \partial_{\bar\zeta}\tilde\varphi\cdot
    \Tr[R_0]\,d^2\zeta}_{=\,0 \;(\text{$R_0$ holomorphic in }\zeta)}
  \;+\;
  \underbrace{-\frac{1}{\pi}\int_\C
    \partial_{\bar\zeta}\tilde\varphi(\zeta)\cdot
    \Tr[R_{\mathrm{Ai}}(\zeta)]\,d^2\zeta}_{=:\,N_{\mathrm{Ai}}(\varphi)}
  \;+\; O(c^{-1/3}).
\end{equation}
The key step: the $c^{2/3}$-term vanishes because $R_0$ is
holomorphic in $\zeta$ (bulk contribution).
The leading term is then:
\begin{equation}\label{eq:NAi}
  N_{\mathrm{Ai}}(\varphi)
  = -\frac{1}{\pi}\int_\C \partial_{\bar\zeta}\tilde\varphi(\zeta)\cdot
  \Tr[(\mathcal{A}-\zeta)^{-1}]\,d^2\zeta
  = \#\{\text{eigenvalues of }\mathcal{A}\text{ in }\mathrm{supp}(\varphi)\}.
\end{equation}

\begin{proposition}[Key reduction --- conditional on \eqref{eq:resolventexp}]
\label{prop:keyreduction}
Assume the resolvent expansion~\eqref{eq:resolventexp} holds in trace norm
for $\zeta$ in a compact set and $c \to \infty$.
Then:
\[
  N_\varphi(c) = N_{\mathrm{Ai}}(\varphi) + O(c^{-1/3}),
\]
where $N_{\mathrm{Ai}}(\varphi)$ counts eigenvalues of the Airy operator
$\mathcal{A} = -d^2/d\xi^2 + \xi$ on $L^2(\R_+)$ in the window $\varphi$.
In particular, if $\varphi$ is adapted to an interval containing
exactly two consecutive zeros of the Airy function $\Ai$,
then $N_\varphi(c) = 2 + O(c^{-1/3})$,
i.e.\ condition~(Q5) holds.
\end{proposition}

\begin{remark}[Why two eigenvalues]
The Airy operator $\mathcal{A}$ on $L^2(\R_+)$ with Dirichlet BC
has eigenvalues at the zeros $-a_k$ of $\Ai$ ($a_1 < a_2 < \cdots$).
Consecutive eigenvalue spacings satisfy
$a_{k+1} - a_k \sim \pi / \sqrt{a_k} \to 0$ as $k\to\infty$,
but for the \emph{first few} zeros the spacings are of order $1$
(in the $c^{-1/3}$-scaled variable).
A window of width $\sim c^{-1/3}$ in the original variable
corresponds to a window of width $\sim 1$ in the $\zeta$-variable,
which generically contains exactly one Airy eigenvalue.
The count of $2$ requires the window to straddle two Airy zeros;
this is the content of matching the PSWF transition index $n^*$
to the correct Airy eigenvalue.
\end{remark}

%% ============================================================
\section{The Airy Operator as Limit}
\label{sec:airyop}
%% ============================================================

\subsection{The prolate spheroidal wave equation near the turning point}

PSWFs $\psi_n^{(c)}$ satisfy the PSWE:
\begin{equation}\label{eq:PSWE}
  \frac{d}{dx}\bigl[(1-x^2)\psi'\bigr] + \bigl[\chi_n - c^2 x^2\bigr]\psi = 0,
  \quad x \in [-1,1],
\end{equation}
where $\chi_n = \chi_n(c)$ is the $n$-th prolate eigenvalue.
In the transition regime $n \sim 2c/\pi$, $x$ near $\pm 1$,
the substitution $x = 1 - c^{-2/3}\xi$ and
$\chi_n = c^2 - 2c^{4/3}(n - 2c/\pi)/\pi + \cdots$ transforms~\eqref{eq:PSWE} into
\begin{equation}\label{eq:Airyeq}
  \frac{d^2 u}{d\xi^2} = \xi \cdot u + o(1),
  \quad \xi \in [0,\infty),
\end{equation}
which is precisely the Airy equation.
The Airy operator $\mathcal{A} = -d^2/d\xi^2 + \xi$ therefore
emerges as the \emph{leading-order limit} of the PSWE in the
double-scaling regime.

\begin{remark}[Known results]
The Airy approximation for PSWFs near the turning point is
known heuristically (Osipov--Rokhlin--Xiao~\cite{Osipov2013})
and asymptotically at the level of individual functions.
What is \emph{not} known is the trace-norm convergence of the
corresponding resolvent, which is the analytic core of this paper.
\end{remark}

\subsection{The Airy operator and its resolvent}

\begin{definition}[Airy operator]
$\mathcal{A} := -d^2/d\xi^2 + \xi$ on $L^2(\R_+)$
with domain $H^2(\R_+) \cap H^1_0(\R_+)$.
\end{definition}

$\mathcal{A}$ is self-adjoint and has discrete spectrum
$\{-a_k : k \geq 1\}$ where $0 < a_1 < a_2 < \cdots$ are the
positive zeros of $\Ai(-\cdot)$.
Eigenfunctions: $u_k(\xi) = c_k \Ai(\xi - a_k)$.
Resolvent for $\zeta \notin \{-a_k\}$:
\begin{equation}\label{eq:Airyres}
  (\mathcal{A}-\zeta)^{-1}u(\xi)
  = \int_0^\infty G_{\mathcal{A}}(\xi,\eta;\zeta)\,u(\eta)\,d\eta,
\end{equation}
where $G_{\mathcal{A}}$ is the Green's function of the Airy equation
with Dirichlet BC, explicitly:
\begin{equation}\label{eq:AiryGreen}
  G_{\mathcal{A}}(\xi,\eta;\zeta)
  = \frac{\pi\,\Ai(\xi_>-\zeta)\,\Bi(\xi_<-\zeta)
         - \pi\,\Bi(\xi_>-\zeta)\,\Ai(\xi_<-\zeta)}
         {W[\Ai,\Bi](\cdot-\zeta)},
\end{equation}
with $\xi_< = \min(\xi,\eta)$, $\xi_> = \max(\xi,\eta)$.

\begin{remark}[Trace of Airy resolvent]
$\Tr[(\mathcal{A}-\zeta)^{-1}] = \sum_k \frac{1}{-a_k - \zeta}$
has poles exactly at the Airy eigenvalues $\{-a_k\}$
and encodes the eigenvalue counting via the residue theorem:
\[
  \#\{k : -a_k \in U\}
  = \frac{1}{2\pi i}\oint_{\partial U} \Tr[(\mathcal{A}-\zeta)^{-1}]\,d\zeta
\]
for any contour $\partial U$ avoiding $\{-a_k\}$.
This is the mechanism by which resolvent control yields counting control.
\end{remark}

%% ============================================================
\section{The Resolvent Convergence Programme}
\label{sec:programme}
%% ============================================================

\subsection{Steps required}

To establish Proposition~\ref{prop:keyreduction} rigorously,
the following chain of statements is required:

\begin{enumerate}[label=\textbf{(S\arabic*)},leftmargin=2.5em]
\item\label{S1} \textbf{Scaled kernel convergence.}
  In the double-scaling~\eqref{eq:scaling}, the kernel of $\Kc$ satisfies
  \[
    c^{-2/3} K_c(1-c^{-2/3}\xi,\; 1-c^{-2/3}\eta)
    \xrightarrow{c\to\infty} K_{\mathcal{A}}(\xi,\eta)
  \]
  in $L^2(\R_+ \times \R_+)$ or in an appropriate trace-class sense,
  where $K_{\mathcal{A}}$ is the spectral projection kernel of $\mathcal{A}$
  below some threshold.

\item\label{S2} \textbf{Scaled resolvent convergence.}
  The scaled resolvent
  \[
    c^{1/3}(\Kc - z_n^* - c^{-1/3}\zeta)^{-1}\big|_{\text{scaled domain}}
    \to (\mathcal{A} - \zeta)^{-1}
  \]
  in trace-norm on the scaled domain,
  uniformly for $\zeta$ in compact subsets of $\C\setminus\R$.

\item\label{S3} \textbf{Trace-norm bound for the tail.}
  The contribution from the bulk (eigenvalues far from $W_c$)
  to $\Tr[\varphi_c(\Kc)]$ is $O(c^{-1/3})$.

\item\label{S4} \textbf{Window calibration.}
  The window $W_c(\zeta_0)$ in the $\zeta$-variable corresponds to
  an interval $[-\zeta_0, \zeta_0]$ containing exactly two zeros
  of $\Ai(-\cdot)$ for the appropriate $\zeta_0 > 0$
  (i.e.\ $[-\zeta_0,\zeta_0] \supset \{-a_1,-a_2\}$
  but $-a_3 \notin [-\zeta_0,\zeta_0]$).
\end{enumerate}

\medskip
\begin{center}
\renewcommand{\arraystretch}{1.3}
\begin{tabular}{llp{4cm}l}
\toprule
\textbf{Step} & \textbf{Content} & \textbf{Key tool} & \textbf{Status}\\
\midrule
\ref{S1} & Kernel convergence & PSWE turning-point asymptotics & open \\
\ref{S2} & Resolvent convergence & Trace-norm perturbation theory & open \\
\ref{S3} & Tail control & Weyl remainder, Gram stability & open \\
\ref{S4} & Window calibration & Zeros of $\Ai$: known explicitly & \checkmark\ (literature) \\
\bottomrule
\end{tabular}
\end{center}

\subsection{The leading analytic difficulty: step (S2)}

Step~\ref{S2} is the hardest. The core obstacle is:

\begin{problem}[Trace-norm resolvent convergence]
\label{prob:tracenorm}
Let $\Kc^{(\mathrm{sc})}$ denote the scaled version of $\Kc$
on $L^2([0,c^{2/3}])$ (after the substitution~\eqref{eq:scaling}).
Prove:
\begin{equation}\label{eq:tracenormconvergence}
  \bigl\|(\Kc^{(\mathrm{sc})} - \zeta)^{-1} - (\mathcal{A}-\zeta)^{-1}\bigr\|_{\mathrm{tr}}
  = O(c^{-1/3})
\end{equation}
uniformly for $\zeta$ in compact subsets of $\C\setminus\R$.
\end{problem}

\begin{remark}[Why trace-norm, not operator-norm]
Operator-norm convergence of the resolvent would follow from
strong resolvent convergence of $\Kc^{(\mathrm{sc})}$ to $\mathcal{A}$,
which may be accessible via form convergence (PSWE $\to$ Airy ODE).
However, for the trace formula~\eqref{eq:HS} to apply,
one needs trace-norm (= nuclear norm) convergence.
This requires $\|(\Kc^{(\mathrm{sc})} - \mathcal{A})\cdot(\mathcal{A}-\zeta)^{-2}\|_{\mathrm{tr}} = O(c^{-1/3})$,
which in turn requires: the difference of the two kernels
has a controlled $L^2$-trace,
depending on the rate of kernel convergence in step~\ref{S1}.
\end{remark}

\subsection{The upper vs.\ lower counting bound}

Proposition~\ref{prop:keyreduction} splits (Q5) into two separate inequalities:

\begin{problem}[(Q5-upper): No spectral crowding]
\label{prob:upper}
Prove $N_\varphi(c) \leq 2 + O(c^{-1/3})$
(at most two eigenvalues in $W_c$).
This is the harder direction: it requires
the \emph{full} resolvent expansion~\eqref{eq:resolventexp}
to exclude extra eigenvalues crowding the window.
\end{problem}

\begin{problem}[(Q5-lower): No eigenvalue escape]
\label{prob:lower}
Prove $N_\varphi(c) \geq 2 - O(c^{-1/3})$
(at least two eigenvalues in $W_c$).
This is the easier direction: quasimode injection.
If two approximate eigenvectors $f_n^{(c)}, f_{n+1}^{(c)}$
with centre $\mu_n, \mu_{n+1} \in W_c$ and residual $\varepsilon = o(c^{-1/3})$ exist,
then $\Tr[\varphi_c(\Kc)] \geq 2 - O(\varepsilon/c^{-1/3})$.
\end{problem}

%% ============================================================
\section{Connection to Fredholm Determinants}
\label{sec:fredholm}
%% ============================================================

An equivalent formulation of the resolvent programme uses
the Fredholm determinant $D_c(\lambda) := \det(I - \lambda \Kc)$.

\begin{remark}[Fredholm determinant and local zeros]
For $\mathrm{Im}(z) \neq 0$:
\[
  \frac{d}{dz}\log D_c(1/z) = \frac{1}{z^2}\Tr[(\Kc - z)^{-1}].
\]
Local eigenvalue counting in $W_c$ corresponds to counting zeros of
$\lambda \mapsto D_c(1/\lambda)$ in a $c^{-1/3}$-neighbourhood.
In the double-scaling limit, the prediction is:
\begin{equation}\label{eq:Dscaling}
  \log D_c(1/(z_n^* + c^{-1/3}\zeta))
  \sim \log D_{\mathcal{A}}(\zeta) + O(c^{-1/3} \log c),
\end{equation}
where $D_{\mathcal{A}}(\zeta) := \det(\mathcal{A} - \zeta)$
(suitably regularised Fredholm determinant of the Airy operator).
The zeros of $D_{\mathcal{A}}$ are precisely the Airy eigenvalues $\{-a_k\}$,
and a Rouché-type argument would then identify local zeros of
$D_c$ with those of $D_{\mathcal{A}}$.
This is the cleanest route to (Q5) if~\eqref{eq:Dscaling} can be established.
\end{remark}

%% ============================================================
\section{Programme State and Open Problems}
\label{sec:state}
%% ============================================================

\begin{center}
\renewcommand{\arraystretch}{1.35}
\begin{tabular}{p{4.0cm} p{3.2cm} p{2.5cm} l}
\toprule
\textbf{Claim} & \textbf{Content} & \textbf{Tool} & \textbf{Status}\\
\midrule
Kernel convergence (S1)
  & Sinc $\to$ Airy in $L^2$
  & PSWE asymptotics & open \\
Resolvent convergence (S2)
  & Trace-norm, $O(c^{-1/3})$
  & Perturbation theory & open (hardest) \\
Tail control (S3)
  & Bulk contrib.\ $O(c^{-1/3})$
  & Weyl + Gram & open \\
Window calibration (S4)
  & 2 Airy zeros in $W_c$
  & Zeros of $\Ai$
  & \checkmark\ known \\
Prop.~\ref{prop:keyreduction}
  & $N_\varphi = 2+O(c^{-1/3})$
  & (S1)--(S4)
  & conditional \\
(Q5-upper) (Prob.~\ref{prob:upper})
  & No crowding
  & Full expansion & open (hard) \\
(Q5-lower) (Prob.~\ref{prob:lower})
  & No escape
  & Quasimode injection & open (easier) \\
Fredholm route (Rem.~4.1)
  & Rouché for $D_c$
  & Det.\ asymptotics & open \\
\bottomrule
\end{tabular}
\end{center}

\medskip\noindent
\textbf{Logical chain (Paper XIV goal):}
\[
  \text{(S1) + (S2) + (S3) + (S4)}
  \;\Longrightarrow\;
  N_\varphi(c) = 2 + O(c^{-1/3})
  \;\Longrightarrow\;
  \text{(Q5)}
  \;\Longrightarrow\;
  \text{Gap-S (via Theorem~B, Paper~XIII).}
\]

\begin{thebibliography}{99}
\bibitem{paper1} \textit{Paper~I: Gram Coercivity}, \texttt{paper1.tex}, 2026.
\bibitem{paper13} \textit{Paper~XIII: Obstructions to Gap-S},
  \texttt{paper13\_gap\_s.tex}, 2026.
\bibitem{Kato1966} T.~Kato, \textit{Perturbation Theory for Linear Operators},
  Springer, 1966.
\bibitem{Osipov2013} A.~Osipov, V.~Rokhlin, H.~Xiao,
  \textit{Prolate Spheroidal Wave Functions of Order Zero}, Springer, 2013.
\bibitem{Widom1964} H.~Widom, \textit{Asymptotic behavior of the eigenvalues
  of certain integral equations},
  Trans.\ Amer.\ Math.\ Soc.\ \textbf{109} (1964), 278--295.
\bibitem{TracyWidom1994} C.~Tracy, H.~Widom, \textit{Level-spacing
  distributions and the Airy kernel},
  Comm.\ Math.\ Phys.\ \textbf{159} (1994), 151--174.
\bibitem{HelSjo1989} B.~Helffer, J.~Sj\"ostrand,
  \textit{\'Equation de Schr\"odinger avec champ magn\'etique},
  Lecture Notes in Physics \textbf{345}, Springer, 1989.
\bibitem{Simon2005} B.~Simon, \textit{Trace Ideals and Their Applications},
  2nd ed., AMS, 2005.
\end{thebibliography}

\end{document}
